% Generated draft: requires mathematical and source review.
\documentclass[11pt]{article}
\usepackage[margin=25mm]{geometry}
\usepackage{fontspec}
\setmainfont{Latin Modern Roman}
\usepackage{amsmath,amssymb,mathtools,unicode-math}
\setmathfont{Latin Modern Math}
\usepackage{enumitem,needspace,xurl,hyperref}
\setlength{\parindent}{0pt}
\setlength{\parskip}{5pt}
\setlength{\emergencystretch}{3em}
\newcommand{\N}{\mathbb{N}}
\newcommand{\Z}{\mathbb{Z}}
\newcommand{\Q}{\mathbb{Q}}
\newcommand{\R}{\mathbb{R}}
\newcommand{\C}{\mathbb{C}}
\newcommand{\F}{\mathbb{F}}
\newcommand{\A}{\mathbb{A}}
\newcommand{\PP}{\mathbb{P}}
\newcommand{\E}{\mathbb{E}}
\newcommand{\eps}{\varepsilon}
\newcommand{\dd}{\,\mathrm{d}}
\begin{document}
\section*{UnsolvedMath ID: 30002690}
\textbf{Source problem code:} catalogue entry
\subsection{Definitions and mathematical statement}
A Riemannian metric \(g\) on a smooth manifold \(M\) is Einstein if
\[
\operatorname{Ric}_g=\lambda g
\qquad\text{for some constant }\lambda\in\R.
\]
A homogeneous Riemannian manifold \((M,g)\) is naturally reductive if it admits a presentation \(M=L/K\), where \(L\) acts transitively by isometries, together with an \(\operatorname{Ad}(K)\)-invariant decomposition
\[
\mathfrak l=\mathfrak k\oplus\mathfrak m
\]
and the inner product on \(\mathfrak m\cong T_{eK}M\) induced by \(g\), such that
\[
\langle [X,Y]_{\mathfrak m},Z\rangle
+\langle Y,[X,Z]_{\mathfrak m}\rangle=0
\qquad\text{for all }X,Y,Z\in\mathfrak m.
\]
Here \([X,Y]_{\mathfrak m}\) denotes projection of the Lie bracket onto \(\mathfrak m\). Under this convention, a homogeneous metric is non-naturally reductive if no such transitive isometric presentation exists. Failure of the displayed identity for just one presentation does not establish this property.

\textbf{Source problem.} Find further non-naturally reductive Einstein metrics \(g\) on compact simple Lie groups \(G\), using homogeneous spaces \(G/H\). A quotient \(G/H\) is understood to be a smooth homogeneous space, so \(H\) is a closed subgroup of \(G\). The metric being sought is on \(G\), whereas \(G/H\) is to enter its construction.

\textbf{Unresolved scope.} The source gives an example-construction task, without specifying a particular group \(G\), subgroup \(H\), construction, or reference collection relative to which examples must be ``further.'' It does not assert existence for every compact simple \(G\), and it does not explicitly request an exhaustive classification. In the usual connected Lie-group convention, simplicity means that the real Lie algebra \(\mathfrak g\) is simple; the record does not specify connectedness, simple connectedness, or the allowed global forms. The context suggests homogeneous metrics on \(G\), commonly left-invariant metrics, where left invariance means \(L_a^*g=g\) for every \(a\in G\), with \(L_a(x)=ax\). However, the source does not explicitly impose left invariance or specify the transitive isometry group. It also does not state whether natural reductivity is to be tested against all transitive isometry groups, as defined above, or against a restricted family of presentations. These choices require source or mathematical review before a more specialized quantified formulation can be asserted.
\subsection{Short English statement}
The task is to construct additional Einstein metrics on compact simple Lie groups that are non-naturally reductive, using quotient spaces \(G/H\) in the construction. The source leaves the groups, subgroups, meaning of ``additional,'' and precise invariance conventions unspecified. Its supplied partially solved status is dated curation, not an independently verified conclusion.
\subsection{Sources}
\begin{itemize}
\item Local UnsolvedMath registry, permanent ID \(30002690\), source problem code \texttt{catalogue entry}, title: \emph{Non-Naturally Reductive Einstein Metrics on Compact Lie Groups}. The supplied record belongs to the Oberwolfach Reports Open Problems collection. No individual proposer is supplied; the proposed year is recorded as 2014.
\item Supplied source citation: \emph{Mini-Workshop: Einstein Metrics, Ricci Solitons and Ricci Flow under Symmetry Assumptions} (2015), Oberwolfach Reports, pp.\ 2529--2568, DOI \texttt{10.4171/owr/2014/45}. \url{https://doi.org/10.4171/owr/2014/45}. The literature-source list labels this reference \emph{Oberwolfach Report 45/2014}. The supplied 2015 citation year and 2014 report designation are retained without reconciliation.
\item Supplied literature reference: \emph{Flag Manifolds with Five Isotropy Summands}. \url{https://doi.org/10.48550/arXiv.1207.2897}. No authors or explanation of its precise relevance are supplied.
\item Literature assessment checked at \texttt{2026-08-21}: the supplied curation labels the problem partially solved and reports subsequent construction and classification of additional non-K\"ahler and non-naturally reductive homogeneous Einstein metrics in several families, while reporting that finding and classifying all such metrics on compact simple groups remains open. This is supplied dated curation, not independently verified literature assessment. Its non-K\"ahler qualification and broader classification discussion are background claims, not additional hypotheses or requirements of the stated construction problem.
\item The registry labels the statement \texttt{corrected\_verified} and reports a manual semantic, typographic, and LaTeX review of a meaning-preserving cleanup. This supplied editorial verification does not establish mathematical correctness or independently verify the literature status. No external source was consulted in preparing this draft.
\end{itemize}
\end{document}
